% ============================================================
%  THE ORTYX PROTOCOL
%  Part II — Tension
%  Chapter 9: When Everything Becomes Evidence
% ============================================================

\chapter{When Everything Becomes Evidence}
\label{ch:semantic-universality}

\chapepigraph{A framework that can explain the presence of a
phenomenon and also explain its absence, using the same
conceptual resources, has not explained either. It has built
a vocabulary in which both outcomes are describable. That is
not the same thing.}{Flyxion}

\noindent
The previous two chapters examined failure modes that operate
at the level of individual claims: the use of notation that
performs constraint without providing it (Decorative Formalism),
and the embedding of conclusions in definitional choices
(Definitional Closure). This chapter examines a failure mode
that operates at the level of the framework as a whole:
the expansion of a framework's abstraction level to the point
where it ceases to be possible to specify what would count
as evidence against its central claims.

This failure mode --- Semantic Universality Collapse
($\FailGeo{3}$) --- is in some ways the most insidious of
the four, because it is often reached by a series of moves
that each appear individually reasonable. A framework
extends its explanatory scope from its original domain
to an adjacent domain; this extension appears successful;
the framework extends again; and so on. At each step,
the framework's vocabulary broadens slightly to accommodate
the new domain. The broadening feels like generalization,
which is a genuine scientific virtue. But at some point
the vocabulary has broadened enough that it can accommodate
\textit{any} observation in any domain, and at that point
generalization has become vacuity.

\section{The Trajectory of Temporal Coherence}
\label{sec:coherence-trajectory}

Tracing the concept of temporal coherence through the Phoenix
corpus reveals a trajectory that exemplifies this dynamic
precisely. The concept begins with a precise, measurable
definition in the domain of audio signal processing: temporal
coherence is low jitter at the peak level of a waveform.
This definition is specific, operationalized, and allows
clear distinctions between signals that have it and signals
that do not.

As the corpus extends the concept to auditory perception,
the definition broadens: temporal coherence is what the
auditory system uses to group spectral components into
streams. This is still reasonably specific --- it can be
studied in auditory scene analysis experiments --- but
the connection to the engineering definition requires a
hypothesis (that the auditory system is tracking something
like peak-level jitter) that the corpus does not establish.

As the concept extends to memory, the definition broadens
further: temporal coherence is the property of interference
patterns in the \MEMeight{} field that makes memories stable
and retrievable. The connection to the auditory definition
is now analogical rather than derivational.

As the concept extends to AI architecture, it broadens
again: temporal coherence is what distinguishes meaningful
representations from noise in any cognitive system. At
this level, the definition is broad enough that any feature
of a representation that makes it useful could potentially
be described as temporal coherence in some sense.

As the concept extends to consciousness, the definition
becomes: temporal coherence is the binding principle of
unified conscious experience. At this level, the concept
is so broad that it is difficult to specify what a
conscious experience that lacked temporal coherence would
be like --- which means it is difficult to specify what
evidence against the claim would look like.

\begin{auditprop}
\label{ap:semantic-universality}
Semantic Universality Collapse ($\FailGeo{3}$) occurs when
a framework's central concept undergoes successive broadening
across domain extensions until it reaches a level of
abstraction at which it can accommodate any observation
in any domain. At this point, the concept ceases to function
as a constraint on what the framework predicts and begins
functioning as a vocabulary for redescribing whatever is
observed. The collapse is gradual rather than sudden,
and the point at which the concept loses its constraining
force is not typically visible from inside the framework.
\end{auditprop}

\section{The Explanatory Expansion of Jitter}
\label{sec:jitter-expansion}

A concrete illustration of this dynamic is provided by
tracing how the corpus handles observations that initially
appear inconsistent with the temporal coherence hypothesis.

Consider the salience of unexpected events --- the point
raised in Chapter~6. High-jitter events are, by the Marine
metric, low-salience. But sudden, unexpected events (breaking
glass, unexpected modulations) are experientially highly
salient. The temporal coherence framework, applied naively,
predicts that these events should not capture attention;
experience says they do.

The framework's response to this challenge, when it responds
at all, is to expand the temporal coherence concept. The
salience of unexpected events, on this expanded reading,
reflects the auditory system's detection of a coherence
violation: the sudden event is salient precisely because
it disrupts a previously coherent temporal model. This
is now called ``coherence contrast'' rather than
``coherence,'' but the concept of coherence is doing the
work in both cases: coherent signals are salient because
they have coherence; incoherent signals are salient because
they violate coherence.

This expansion is not obviously wrong. Contrast effects
are real, and there is genuine perceptual evidence that
the salience of a stimulus depends partly on how different
it is from the immediately preceding context. But the
expansion has a specific structural consequence: the
framework can now accommodate both coherent salience and
incoherent salience as manifestations of temporal coherence
dynamics. No observation about salience can disconfirm
the temporal coherence hypothesis, because any salient
stimulus can be described either as having coherence or
as violating coherence, and both descriptions invoke the
framework.

\begin{dangerzone}[Universal Absorbency]
A framework has entered the regime of Universal Absorbency
when its central concept has been expanded to accommodate
both the presence and the absence of the phenomenon it
was introduced to explain. The temporal coherence concept,
extended to include both ``salience from coherence'' and
``salience from coherence violation,'' can accommodate
any pattern of attention in any signal. This means the
temporal coherence framework, as extended, does not predict
what patterns of attention will be observed; it provides
a vocabulary for redescribing whatever patterns are
observed after the fact. That is a significant epistemic
change from the original, specific claim.
\end{dangerzone}

\section{The Problem of Consistent Accommodation}
\label{sec:consistent-accommodation}

Semantic Universality Collapse does not require that the
framework explicitly accommodate anomalous observations
by ad hoc modification. It can occur through a subtler
mechanism: the framework's vocabulary is broad enough
from the start that no clear anomaly can arise.

Consider the concept of ``authentic experience'' as used
in the corpus. If authentic experience is associated with
temporally coherent signals, what would an inauthentic
experience be? The corpus implies that inauthentic
experience is what listeners have when listening to
compressed audio --- that spectral fracture produces a
kind of experiential inauthenticity. But what would it
mean for a listener to have an authentic experience of
a compressed recording, or an inauthentic experience of
an uncompressed one?

If authentic experience just is the experience of temporally
coherent signals, then by definition uncompressed audio
produces authentic experience and compressed audio does
not. The question of whether any listener actually finds
their experience of compressed audio inauthentic becomes
irrelevant to whether compressed audio is, by the
framework's lights, inauthentic. The framework has made
an empirical-sounding claim (compressed audio produces
inauthentic experience) that is actually definitional
(compressed audio is, by definition, what inauthentic
audio is).

\section{Distinguishing Genuine Generalization}
\label{sec:genuine-generalization}

Semantic Universality Collapse must be distinguished from
genuine scientific generalization, which involves
real extension of explanatory power. The distinction
is not always easy, but it is principled.

Genuine generalization maintains the falsifiability of
the extended claims. When classical mechanics was extended
from terrestrial to celestial mechanics, the extension
generated new predictions --- about the orbits of planets,
about the precession of equinoxes, about the return of
comets --- that could in principle have been false.
The extension was successful because the predictions were
largely confirmed, not because the framework was defined
in a way that made failure impossible.

Atmospheric generalization, by contrast, applies a
framework's vocabulary to new domains without generating
new falsifiable predictions. The application organizes
the new domain within the framework's conceptual structure
and makes claims that sound predictive but that are
actually redescriptions: ``this phenomenon exhibits
temporal coherence dynamics'' where any phenomenon would.

The test is always the same: what would the domain
look like if the extended framework were false? If the
answer is ``exactly as it looks now, described in different
terms,'' then the generalization is atmospheric rather
than genuine.

For the temporal coherence framework, the test at each
level of extension is:

At the audio engineering level: what signals would be
classified as meaningful by the framework but would fail
to be experienced as meaningful by listeners? If such
signals are possible in principle, the framework is making
genuine predictions. If the framework would accommodate
any listener response by adjusting its parameters, it
is not.

At the memory level: what memory phenomena would occur
in a system with low temporal coherence in its field
dynamics? If such phenomena are possible in principle,
the extended framework is making genuine predictions.
If any memory phenomenon can be redescribed as a
manifestation of field coherence or field incoherence,
it is not.

At the consciousness level: what conscious experiences
would occur in the absence of temporal coherence? If
the framework implies that consciousness is impossible
without temporal coherence, it is making a strong and
potentially falsifiable claim. If it implies that any
conscious experience can be characterized in terms of
some form of coherence dynamics, it is not.

The Phoenix corpus does not address these tests explicitly.
This is not a final verdict; it is an observation about
what remains to be done for the framework to achieve the
status of genuine generalization rather than atmospheric
expansion.

\section{The Epistemic Thermodynamics of Scope Inflation}
\label{sec:epistemic-thermodynamics}

There is a sense in which Semantic Universality Collapse
represents a kind of epistemic thermodynamics. A framework
begins with high information content: it makes specific
predictions that exclude many possible observations.
As it extends its scope by broadening its vocabulary,
its information content decreases: each broadening
accommodates more possible observations. If the broadening
continues without a corresponding increase in the
specificity of the framework's claims, the framework
asymptotically approaches a state of maximum epistemic
entropy: it can accommodate all possible observations
and therefore carries no information about which
observations will actually be made.

This thermodynamic framing is metaphorical --- epistemic
entropy is not the same as thermodynamic entropy --- but
it captures something structurally real. A framework that
has undergone Semantic Universality Collapse has maximized
its accommodation capacity at the cost of its predictive
content. It has traded information for coverage.

The trading is not always a failure. At the level of
computational metaphor, a framework that organizes a wide
range of phenomena within a unified vocabulary may be
providing genuine intellectual value even if it is not
making testable predictions. The value is organizational
and heuristic rather than predictive. The failure occurs
when the framework's practitioners and readers mistake
organizational coverage for predictive content --- when
the appearance of unified explanation is taken as evidence
of a unified mechanism.

\begin{epistemicstatus}{Phenomenological}
The experience of reading a framework that has undergone
Semantic Universality Collapse is phenomenologically
distinctive: everything seems to fit, connections appear
everywhere, the framework feels illuminating, and it
becomes difficult to formulate a clear objection because
the framework's vocabulary is available to redescribe
any objection as a further confirmation. This is not an
argument against the framework; it is a description of
the epistemic experience that warrants suspicion. When
everything seems to fit, the appropriate response is
not satisfaction but the question: what would not fit?
\end{epistemicstatus}

\medskip

Chapter~10 examines the fourth and final failure geometry:
Operational Severance, the condition in which a framework's
formal vocabulary loses its connection to measurement
procedures. This is the failure mode most specific to
the current period of intellectual production, in which
the technical apparatus of data science, information
theory, and computational neuroscience provides a rich
vocabulary of formally dressed claims that can be deployed
in domains where no measurement procedure for their
key quantities has been established.
